\section{优化模型}
% 替换为按以下顺序撰写的模型内容（完成后删除本骨架注释）。
% 建模假设：编号列出对模型有实质影响的条件，并简述理由。
% 符号与数据约定：列出符号、含义、单位和来源。
% 决策变量：逐个说明符号、取值域和含义。
% 目标与约束：写出目标函数并解释业务含义；按业务目的分组约束，给出公式与简要解释。说明模型相对原题的变换。
% 模型规模：汇总变量数、约束行数和整数变量数。
% 最后一小节“完整优化模型”：集中列出完整业务模型或已核实等价的完整表达，包括目标、全部必要约束及索引/量词、所有变量定义域。可引用上文已定义的集合和参数；不得用松弛或分解子问题、或“其余约束同上”代替完整模型。
% 注：约束正确性需要论证时，默认用正文一句话 + 附录数值小算例（附录用 \appendix）；默认不使用命题/证明环境。
